\documentclass{article}
\usepackage{amsmath, amssymb, amsthm}
\usepackage{hyperref}
\usepackage{geometry}
\geometry{margin=1in}

\title{Erd\H{o}s Problem \#1183}
\date{Last updated: 2026-03-07}
\author{Source: erdosproblems.com}

\begin{document}
\maketitle

\noindent\textbf{Status:} OPEN \quad \textbf{No prize}

\noindent\textbf{Tags:} combinatorics, ramsey theory

\medskip

\noindent\textbf{Problem Statement:}

Let $f(n)$ be maximal such that in any $2$-colouring of the subsets of $\{1,\ldots,n\}$ there is always a monochromatic family of at least $f(n)$ sets which is closed under taking unions and intersections. Estimate $f(n)$. Let $F(n)$ be defined similarly, except that we only require the family be closed under taking unions. Estimate $F(n)$. In particular, is it true that $F(n)\geq n^{\omega(n)}$ for some $\omega(n)\to \infty$ as $n\to \infty$, and $F(n)&#60;(1+o(1))^n$?




A problem of Erd\H{o}s and Ulam. It is trivial that $f(n)\geq \frac{n+1}{2}$, since there is a sequence of $n+1$ many nested subsets. This seems to be all that they knew; Erd\H{o}s \cite{Er78} wrote 'we have no plausible conjecture for the true order of magnitude of $f(n)$'. They similarly had no good guess about $F(n)$. In \cite{Er78} Erd\H{o}s wrote that if the colouring is such that all subsets of same size receive the same colour then Howorka had proved that $F(n)&#62;n^{\omega(n)}$ for some $\omega(n)\to \infty$, but gave no reference.




References

[Er78] Erd\H{o}s, Paul, Problems and results in combinatorial analysis and combinatorial number theory . Proceedings of the Ninth Southeastern Conference on Combinatorics, Graph Theory, and Computing (Florida Atlantic Univ., Boca Raton, Fla., 1978) (1978), 29-40.

\noindent\textbf{Related OEIS sequences:} possible


\bigskip
\noindent\small{Source: \url{https://www.erdosproblems.com/1183}}
\end{document}
